
برای عدد صحیح
\(p\geq1\)
فرض کنید
\(x^*\)
نقطهٔ ثابت تابع
\(g\)
است. دنبالهٔ
\(\{x^k\}\)
با شروع از
\(x^0\)
و با رابطهٔ
\[
 x^{k+1}=g(x^k)
\]
تولید شده و به
\(x^*\)
همگراست و تکرار در تعداد متناهی گام در
\(x^*\)
متوقف نمی‌شود؛ یعنی برای همه‌ی (k)های به اندازه‌ی کافی بزرگ،
\(x^k\neq x^*\).
فرض کنید
	\[g^{(j)}(x^*)=0,\qquad j=1,\ldots,p-1.\]
	(برای \(p=1\)، این فهرست از شرط‌ها تهی است.)
	و \(g^{(p)}(x)\)
 	به ازای همه 
 	\(x\)
 	در همسایگی
 	\(x^*\)
 	به شعاع
 	\(\epsilon\)
 	یعنی در بازه
 	\((x^* -\epsilon, x^*+\epsilon)\)
	پیوسته است و
	\(g^{(p)}(x^*)\neq0\).
 در این صورت نشان دهید مرتبهٔ همگرایی دنباله
  دقیقاً \(p\)
 است. اگر شرط آخر حذف شود، نشان دهید مرتبهٔ همگرایی دست‌کم \(p\) است.
 \\
\textbf{یادآوری: تعریف همگرایی  مرتبه 
	\(p\) ام
	:}
فرض کنید دنباله
\(\{x^k\}\)
به 
\(x^*\)
همگراست . اگر به ازای مقدار ثابت 
\(\bar{c}\)
و عدد صحیح
\(p\geq 1\)
داشته باشیم:
\[\|x^{k+1} - x^*\| \leq \bar{c}\|x^k - x^*\|^p\]
آنگاه همگرایی دنباله‌ی
\(\{x^k\}\)
به 
\(x^*\)
دست‌کم از مرتبهٔ
\(p\)
است. برای آن‌که مرتبه دقیقاً
\(p\)
باشد، معیار استاندارد این است که
\[
\lim_{k\to\infty}
\frac{\|x^{k+1}-x^*\|}{\|x^k-x^*\|^p}
=c,
\qquad 0<c<\infty.
\]
\\
 \textbf{راهنمایی:}   استفاده از قضیه تیلور  و پیوستگی تابع مشتق مرتبه
\(p\)
ام در همسایگی  
\(x^*\).

\begin{answer}
	
	\textbf{پاسخ.}
	
\end{answer}
